\textbf{Phần 1. Trắc nghiệm lựa chọn}
\setcounterpageref{ex}{0}

\begin{ex} %[2D8V42]
Trong phép chia $18 : 2 = 9$, số $18$ được gọi là
\choice
{Số chia}
{\True Số bị chia}
{Thương}
{Số hạng}
\loigiai{
18 là số đứng trước dấu chia nên được gọi là Số bị chia.
\textbf{Đáp án của câu hỏi là B}
}
\end{ex}

\begin{ex} %[2D8V42]
Biết số bị chia là $35$, số chia là $5$. Thương là
\choice
{\True $7$}
{$6$}
{$8$}
{$9$}
\loigiai{
Thương $= 35 : 5 = 7$.
\textbf{Đáp án của câu hỏi là A}
}
\end{ex}

\begin{ex} %[2D8V42]
Tìm số bị chia biết số chia là $2$ và thương là $8$.
\choice
{$14$}
{$18$}
{\True $16$}
{$12$}
\loigiai{
Số bị chia = Thương $\times$ Số chia = $8 \times 2 = 16$.
\textbf{Đáp án của câu hỏi là C}
}
\end{ex}

\begin{ex} %[2D8V42]
Tìm số chia biết số bị chia là $20$ và thương là $4$.
\choice
{$4$}
{$6$}
{$3$}
{\True $5$}
\loigiai{
Số chia = Số bị chia : Thương = $20 : 4 = 5$.
\textbf{Đáp án của câu hỏi là D}
}
\end{ex}

\begin{ex} %[2D8V42]
Trong phép chia có số bị chia bằng số chia (khác 0), thương luôn bằng
\choice
{$0$}
{\True $1$}
{Chính số đó}
{Không xác định}
\loigiai{
Một số khác 0 chia cho chính nó luôn có thương bằng 1.
\textbf{Đáp án của câu hỏi là B}
}
\end{ex}

\begin{ex} %[2D8V42]
Thương của số lớn nhất có một chữ số và số $1$ là
\choice
{\True $9$}
{$1$}
{$8$}
{$10$}
\loigiai{
Số lớn nhất có 1 chữ số là 9. Thương: $9 : 1 = 9$.
\textbf{Đáp án của câu hỏi là A}
}
\end{ex}

\begin{ex} %[2D8V42]
Một phép chia có thương bằng $5$, nếu giữ nguyên số chia là $2$ và tăng số bị chia thêm $4$ đơn vị thì thương mới là
\choice
{$6$}
{$8$}
{\True $7$}
{$9$}
\loigiai{
Số bị chia ban đầu: $5 \times 2 = 10$. Số bị chia mới: $10 + 4 = 14$. Thương mới: $14 : 2 = 7$.
\textbf{Đáp án của câu hỏi là C}
}
\end{ex}

\begin{ex} %[2D8V42]
Điền số thích hợp: $\dots : 5 = 8$.
\choice
{$35$}
{$45$}
{$30$}
{\True $40$}
\loigiai{
Số bị chia $= 8 \times 5 = 40$.
\textbf{Đáp án của câu hỏi là D}
}
\end{ex}

\begin{ex} %[2D8V42]
Có $16\text{ L}$ nước mắm chia đều vào các can, mỗi can $2\text{ L}$. Số can nước mắm là
\choice
{$7$ can}
{\True $8$ can}
{$9$ can}
{$6$ can}
\loigiai{
$16 : 2 = 8$ (can).
\textbf{Đáp án của câu hỏi là B}
}
\end{ex}

\begin{ex} %[2D8V42]
Trong một phép chia, số bị chia là số chẵn nhỏ nhất có hai chữ số, số chia là $2$. Thương là
\choice
{\True $5$}
{$4$}
{$6$}
{$7$}
\loigiai{
Số chẵn nhỏ nhất có 2 chữ số là 10. Thương: $10 : 2 = 5$.
\textbf{Đáp án của câu hỏi là A}
}
\end{ex}

\begin{ex} %[2D8V42]
Tìm $y$ biết: $45 : y = 5$.
\choice
{$8$}
{$7$}
{\True $9$}
{$6$}
\loigiai{
$y = 45 : 5 = 9$.
\textbf{Đáp án của câu hỏi là C}
}
\end{ex}

\begin{ex} %[2D8V42]
Một lớp học có $35$ học sinh được chia đều thành các tổ, mỗi tổ có $5$ bạn. Hỏi có bao nhiêu tổ?
\choice
{$6$ tổ}
{$8$ tổ}
{$5$ tổ}
{\True $7$ tổ}
\loigiai{
$35 : 5 = 7$ (tổ).
\textbf{Đáp án của câu hỏi là D}
}
\end{ex}

\textbf{Phần 2. Trắc nghiệm đúng sai}
\setcounterpageref{ex}{0}

\begin{ex} %[2D8V42]
Xét tính đúng sai của các nhận định về thành phần phép chia:
\choiceTF
{\True Công thức: Số bị chia = Thương x Số chia.}
{\True Công thức: Số chia = Số bị chia : Thương (với thương khác 0).}
{Trong phép chia, số bị chia luôn phải nhỏ hơn số chia.}
{Thương của 0 chia cho bất kỳ số nào khác 0 luôn bằng chính số đó.}
\loigiai{
\begin{itemize}
    \item Ý a) Đúng: Đúng.
    \item Ý b) Đúng: Đúng.
    \item Ý c) Sai: Số bị chia trong toán lớp 2 luôn lớn hơn hoặc bằng số chia.
    \item Ý d) Sai: $0 : a = 0$ (bằng 0 chứ không phải bằng chính số đó).
\end{itemize}
\begin{itemize}
    \item a) \textbf{Đúng}.
    \item b) \textbf{Đúng}.
    \item c) \textbf{Sai}.
    \item d) \textbf{Sai}.
\end{itemize}
}
\end{ex}

\begin{ex} %[2D8V42]
Cho phép chia 40 : 5 = 8:
\choiceTF
{\True Số bị chia là 40, số chia là 5, thương là 8.}
{Nếu gấp số chia lên 2 lần thành 10 thì thương tăng gấp đôi.}
{\True Tích của số chia và thương là 5 x 8 = 40.}
{\True Thương 8 là một số chẵn.}
\loigiai{
\begin{itemize}
    \item Ý a) Đúng: Đúng.
    \item Ý b) Sai: Khi số chia tăng thì thương giảm: $40 : 10 = 4$ (giảm đi một nửa).
    \item Ý c) Đúng: Đúng.
    \item Ý d) Đúng: Đúng.
\end{itemize}
\begin{itemize}
    \item a) \textbf{Đúng}.
    \item b) \textbf{Sai}.
    \item c) \textbf{Đúng}.
    \item d) \textbf{Đúng}.
\end{itemize}
}
\end{ex}

\textbf{Phần 3. Trả lời ngắn}
\setcounterpageref{ex}{0}

\begin{ex} %[2D8V42]
Tìm số bị chia biết số chia là 5 và thương là 7.
\shortans{35}
\loigiai{
\textbf{Đáp số:} $35$\\
$7 \times 5 = 35$.
}
\end{ex}

\begin{ex} %[2D8V42]
Tìm số chia biết số bị chia là 16 và thương là 2.
\shortans{8}
\loigiai{
\textbf{Đáp số:} $8$\\
$16 : 2 = 8$.
}
\end{ex}

\begin{ex} %[2D8V42]
Tính thương của 14 và 2.
\shortans{7}
\loigiai{
\textbf{Đáp số:} $7$\\
$14 : 2 = 7$.
}
\end{ex}

\begin{ex} %[2D8V42]
Có 30 bông hoa cắm đều vào các bình, mỗi bình 5 bông hoa. Hỏi cắm được bao nhiêu bình hoa?
\shortans{6}
\loigiai{
\textbf{Đáp số:} $6$\\
$30 : 5 = 6$ bình.
}
\end{ex}

\begin{ex} %[2D8V42]
Tìm một số biết rằng lấy số đó nhân với 5 thì được 45.
\shortans{9}
\loigiai{
\textbf{Đáp số:} $9$\\
$45 : 5 = 9$.
}
\end{ex}

\begin{ex} %[2D8V42]
Một đàn gà có 18 cái chân. Hỏi đàn gà đó có bao nhiêu con gà?
\shortans{9}
\loigiai{
\textbf{Đáp số:} $9$\\
$18 : 2 = 9$ con gà.
}
\end{ex}
